\documentclass[10pt,a4paper]{amsart}
\usepackage{amsmath,amssymb,amsthm,graphicx}
% Retained minimal preamble for the Stacks Project chapter more-algebra.tex (More on Algebra).
% The chapter file pinned-original-source.tex carries only \input{preamble} on its first line, so no
% preamble text lives inside the chapter itself.  The Stacks Project uses one single shared preamble for
% every chapter; the snapshot stages it identically as sources/stacks-chow/preamble.tex and
% sources/stacks-etale-cohomology/preamble.tex.  The block below is copied mechanically, line for line,
% from sources/stacks-chow/preamble.tex lines 14-18 (the Xy-pic setup) and 154-217 (the native theorem
% declarations together with the mathematical macros).  Nothing was added, removed or redefined; the
% wrapper class and the amsmath/amssymb/amsthm/graphicx base packages are supplied by the assembler.
% Chapter: more-algebra.tex -- More on Algebra.
% For commutative diagrams we use Xy-pic
\usepackage[all]{xy}

% We use 2cell for 2-commutative diagrams.
\xyoption{2cell}

% Theorem environments.
%
\theoremstyle{plain}
\newtheorem{theorem}[subsection]{Theorem}
\newtheorem{proposition}[subsection]{Proposition}
\newtheorem{lemma}[subsection]{Lemma}

\theoremstyle{definition}
\newtheorem{definition}[subsection]{Definition}
\newtheorem{example}[subsection]{Example}
\newtheorem{exercise}[subsection]{Exercise}
\newtheorem{situation}[subsection]{Situation}

\theoremstyle{remark}
\newtheorem{remark}[subsection]{Remark}
\newtheorem{remarks}[subsection]{Remarks}

\numberwithin{equation}{subsection}

% Macros
%
\def\lim{\mathop{\mathrm{lim}}\nolimits}
\def\colim{\mathop{\mathrm{colim}}\nolimits}
\def\Spec{\mathop{\mathrm{Spec}}}
\def\Hom{\mathop{\mathrm{Hom}}\nolimits}
\def\Ext{\mathop{\mathrm{Ext}}\nolimits}
\def\SheafHom{\mathop{\mathcal{H}\!\mathit{om}}\nolimits}
\def\SheafExt{\mathop{\mathcal{E}\!\mathit{xt}}\nolimits}
\def\Sch{\mathit{Sch}}
\def\Mor{\mathop{\mathrm{Mor}}\nolimits}
\def\Ob{\mathop{\mathrm{Ob}}\nolimits}
\def\Sh{\mathop{\mathit{Sh}}\nolimits}
\def\NL{\mathop{N\!L}\nolimits}
\def\CH{\mathop{\mathrm{CH}}\nolimits}
\def\proetale{{pro\text{-}\acute{e}tale}}
\def\etale{{\acute{e}tale}}
\def\QCoh{\mathit{QCoh}}
\def\Ker{\mathop{\mathrm{Ker}}}
\def\Im{\mathop{\mathrm{Im}}}
\def\Coker{\mathop{\mathrm{Coker}}}
\def\Coim{\mathop{\mathrm{Coim}}}

% Boxtimes
%
\DeclareMathSymbol{\boxtimes}{\mathbin}{AMSa}{"02}

%
% Macros for moduli stacks/spaces
%
\def\QCohstack{\mathcal{QC}\!\mathit{oh}}
\def\Cohstack{\mathcal{C}\!\mathit{oh}}
\def\Spacesstack{\mathcal{S}\!\mathit{paces}}
\def\Quotfunctor{\mathrm{Quot}}
\def\Hilbfunctor{\mathrm{Hilb}}
\def\Curvesstack{\mathcal{C}\!\mathit{urves}}
\def\Polarizedstack{\mathcal{P}\!\mathit{olarized}}
\def\Complexesstack{\mathcal{C}\!\mathit{omplexes}}
% \Pic is the operator that assigns to X its picard group, usage \Pic(X)
% \Picardstack_{X/B} denotes the Picard stack of X over B
% \Picardfunctor_{X/B} denotes the Picard functor of X over B
\def\Pic{\mathop{\mathrm{Pic}}\nolimits}
\def\Picardstack{\mathcal{P}\!\mathit{ic}}
\def\Picardfunctor{\mathrm{Pic}}
\def\Deformationcategory{\mathcal{D}\!\mathit{ef}}

% --- appended 2026-10-12 (Chief): the project own definitions of the reference/slogan/history
% display environments, copied mechanically from the shared Stacks preamble lines 8-12
% (\usepackage{verbatim} plus the three \newenvironment lines, which the chapter preamble keeps
% verbatim from the project preamble). Nothing added, removed or redefined: these environments
% carry no mathematical content and the project itself typesets them as comments. Without these
% three lines a counted statement whose own text uses them fails with Environment slogan undefined.
\usepackage{verbatim}
\newenvironment{reference}{\comment}{\endcomment}
\newenvironment{slogan}{\comment}{\endcomment}
\newenvironment{history}{\comment}{\endcomment}
\usepackage{cleveref}
\crefname{originallemma}{original lemma}{original lemmas}
\Crefname{originallemma}{Original lemma}{Original lemmas}
\makeatletter
\newlabel{more-algebra-lemma-henselization-local-ring}{{0A03}{source}{}{sourcepointer-1}{}}
\newlabel{more-algebra-lemma-henselization-local-ring@cref}{{[originallemma][0][]0A03}{[1][0][]source}{}{}{}}
\makeatother
\crefname{originallemma}{original lemma}{original lemmas}
\Crefname{originallemma}{Original lemma}{Original lemmas}
\makeatletter
\newlabel{lemma-henselization}{{0A02}{source}{}{sourcepointer-2}{}}
\newlabel{lemma-henselization@cref}{{[originallemma][0][]0A02}{[1][0][]source}{}{}{}}
\makeatother
\crefname{originallemma}{original lemma}{original lemmas}
\Crefname{originallemma}{Original lemma}{Original lemmas}
\makeatletter
\newlabel{lemma-henselization-flat}{{0AGU}{source}{}{sourcepointer-3}{}}
\newlabel{lemma-henselization-flat@cref}{{[originallemma][0][]0AGU}{[1][0][]source}{}{}{}}
\makeatother
\crefname{originallemma}{original lemma}{original lemmas}
\Crefname{originallemma}{Original lemma}{Original lemmas}
\makeatletter
\newlabel{algebra-lemma-henselization}{{04GN}{source}{}{sourcepointer-4}{}}
\newlabel{algebra-lemma-henselization@cref}{{[originallemma][0][]04GN}{[1][0][]source}{}{}{}}
\makeatother
\crefname{originallemma}{original lemma}{original lemmas}
\Crefname{originallemma}{Original lemma}{Original lemmas}
\makeatletter
\newlabel{algebra-lemma-uniqueness-henselian}{{08HT}{source}{}{sourcepointer-5}{}}
\newlabel{algebra-lemma-uniqueness-henselian@cref}{{[originallemma][0][]08HT}{[1][0][]source}{}{}{}}
\makeatother
\title{The henselization functor on local rings agrees with the henselization of the associated pair}
\author{The Stacks Project Authors}
\date{stacks-project:more-algebra.tex:local-snapshot-20261012; GNU FDL 1.2 or later}
\begin{document}
\maketitle
\setcounter{subsection}{1}
\noindent\textit{Editorial note.} All delivered mathematical text is the author's own, copied from the pinned Stacks Project chapter file more-algebra.tex; no statement or proof step was rewritten, reordered or removed.\par\medskip
\noindent\textit{Editorial note.} The chapter file pinned-original-source.tex carries only \textbackslash{}input\{preamble\} on its first line, so no preamble text lives inside the chapter. The retained file sources/stacks-more-algebra/retained-preamble.tex is a mechanical, line-for-line copy of the project's single shared preamble as staged in the snapshot (Xy-pic setup plus theorem declarations and mathematical macros). Nothing was added, removed or redefined, and the excerpt serves only as preamble, never as a span coordinate.\par\medskip
\noindent\textit{Editorial note.} Both counted spans start at the \textbackslash{}begin line of their environment, so the statement and the author's proof are delivered whole; surrounding prose of the section is deliberately outside the counted spans and is carried as standard prerequisites instead.\par\medskip
\noindent\textit{Editorial note.} The citation fields (source\_id, source\_version, authors, source\_url, license.name, source\_title) are identical to those of every other card this lane has produced from this chapter, so the batch selector treats them as one source.\par\medskip
\noindent\textit{Allowed standard prerequisites (editorial).}\begin{itemize}
\item Construction of the henselization of a pair. Lemma \textbackslash{}ref\{lemma-henselization\}, cited at proof lines 3118 and 3129 and 3132-3134. Provides the ring whose maximal ideal and colimit description are studied.
\item The maximal ideal extends to a maximal ideal contained in the Jacobson radical, and the residue field is unchanged. Lemma \textbackslash{}ref\{lemma-henselization-flat\}, cited at proof lines 3119-3120 and 3136. Shows that the henselization of the pair is a local ring with the expected maximal ideal and residue field.
\item The construction of the henselization of a local ring as a colimit of etale algebras. Algebra, Lemma \textbackslash{}ref\{algebra-lemma-henselization\}, cited at proof lines 3127-3128. Allows the two constructions to be compared directly.
\item Uniqueness of the henselian local ring which is a filtered colimit of etale algebras with prescribed residue field. Algebra, Lemma \textbackslash{}ref\{algebra-lemma-uniqueness-henselian\}, cited at proof lines 3137-3138. Gives the canonical isomorphism in the second proof.
\end{itemize}\medskip
\ifdefined\hypertarget\hypertarget{sourcepointer-1}{}\fi
\noindent\textit{Source reference.} Original lemma 0A03: focused-seeds/AGP/sources/stacks-snapshot-20261011/tags.txt, source lines 12621--12621. The lemma's own canonical tag, quoted so that readers can resolve the label inside the delivered text; it points at the counted statement itself and carries no proof content.\par\medskip
\ifdefined\hypertarget\hypertarget{sourcepointer-2}{}\fi
\noindent\textit{Source reference.} Original lemma 0A02: focused-seeds/AGP/sources/stacks-snapshot-20261011/tags.txt, source lines 12620--12620. Navigation pointer to a statement the counted proof refers to but does not contain; nothing load-bearing is delegated to it.\par\medskip
\ifdefined\hypertarget\hypertarget{sourcepointer-3}{}\fi
\noindent\textit{Source reference.} Original lemma 0AGU: focused-seeds/AGP/sources/stacks-snapshot-20261011/tags.txt, source lines 13207--13207. Navigation pointer to a statement the counted proof refers to but does not contain; nothing load-bearing is delegated to it.\par\medskip
\ifdefined\hypertarget\hypertarget{sourcepointer-4}{}\fi
\noindent\textit{Source reference.} Original lemma 04GN: focused-seeds/AGP/sources/stacks-snapshot-20261011/tags.txt, source lines 5839--5839. Navigation pointer to a statement the counted proof refers to but does not contain; nothing load-bearing is delegated to it.\par\medskip
\ifdefined\hypertarget\hypertarget{sourcepointer-5}{}\fi
\noindent\textit{Source reference.} Original lemma 08HT: focused-seeds/AGP/sources/stacks-snapshot-20261011/tags.txt, source lines 10791--10791. Navigation pointer to a statement the counted proof refers to but does not contain; nothing load-bearing is delegated to it.\par\medskip
\begin{lemma}
\label{lemma-henselization-local-ring}
\begin{slogan}
Compatibility henselization of pairs and of local rings.
\end{slogan}
The functor of Lemma \ref{lemma-henselization} associates to a local ring
$(A, \mathfrak m)$ its henselization.
\end{lemma}

\begin{proof}
Let $(A^h, \mathfrak m^h)$ be the henselization of the pair $(A, \mathfrak m)$
constructed in Lemma \ref{lemma-henselization}. Then
$\mathfrak m^h = \mathfrak m A^h$ is a maximal ideal
by Lemma \ref{lemma-henselization-flat} and since it
is contained in the Jacobson radical, we conclude $A^h$ is
local with maximal ideal $\mathfrak m^h$. Having said this
there are two ways to finish the proof.

\medskip\noindent
First proof: observe that the construction in the proof of
Algebra, Lemma \ref{algebra-lemma-henselization} as a colimit
is the same as the colimit used to construct $A^h$ in
Lemma \ref{lemma-henselization}.
Second proof: Both the henselization $A \to S$ and $A \to A^h$ of
Lemma \ref{lemma-henselization} are local ring homomorphisms,
both $S$ and $A^h$ are filtered colimits of \'etale $A$-algebras,
both $S$ and $A^h$ are henselian local rings, and both $S$ and $A^h$
have residue fields equal to $\kappa(\mathfrak m)$ (by
Lemma \ref{lemma-henselization-flat} for the second case).
Hence they are canonically isomorphic by
Algebra, Lemma \ref{algebra-lemma-uniqueness-henselian}.
\end{proof}

\end{document}
